\documentclass[11pt]{article}

\usepackage[margin=1in]{geometry}
\usepackage[T1]{fontenc}
\usepackage{lmodern}

\usepackage{amsmath, amssymb}

\usepackage{setspace}

\usepackage[most]{tcolorbox}
\usepackage{xcolor}

\tcbset{
  enhanced,
  boxrule=0pt,
  frame hidden,
  breakable
}

\usepackage{fancyhdr}
\pagestyle{fancy}
\fancyhf{}
\setlength{\headheight}{52pt}
\renewcommand{\headrulewidth}{0pt}

\newcommand{\Course}{Math 4610 Spring 2026}
\newcommand{\Instructor}{Discussion 5}
\newcommand{\AssignmentType}{Discussion} % or Discussion
\newcommand{\AssignmentNumber}{5}

\newcommand{\Contributors}{
Marks, Stephen\\
}

\fancyhead[L]{\Course\\ \Instructor}
\fancyhead[R]{\Contributors}

\newcommand{\ProblemDivider}[1]{%
  \vspace{1.0em}
  \noindent\textbf{#1}\par
  \vspace{0.35em}
}

\definecolor{problemBack}{RGB}{235,242,250}
\definecolor{problemFrame}{RGB}{70,110,160}
\definecolor{exerciseGray}{gray}{0.96}
\definecolor{exerciseGrayFrame}{gray}{0.45}

\newtcolorbox{problemBox}{
  colback=exerciseGray,
  colframe=exerciseGrayFrame,
  arc=2mm,
  left=8pt,right=8pt,top=8pt,bottom=8pt
}

\newtcolorbox{solutionBox}{
  colback=white,
  colframe=white,
  arc=0mm,
  left=0pt,right=0pt,top=0pt,bottom=0pt
}

\newcommand{\ProblemAnnounce}{\noindent}
\newcommand{\SolutionAnnounce}{\noindent\textbf{Solution.}\par\medskip}

\newcommand{\Exercise}[1]{\textbf{Exercise~#1.}\;}

\newcommand{\R}{\mathbb{R}}
\newcommand{\C}{\mathbb{C}}
\newcommand{\F}{\mathbb{F}}

\begin{document}

\begin{center}
  {\Large \textbf{\AssignmentType~\AssignmentNumber}}
\end{center}
\vspace{0.75em}

\ProblemDivider{1}

\begin{problemBox}
\ProblemAnnounce
\Exercise{5.D.21}
The Fibonacci sequence \(F_{0},F_{1},F_{2},\ldots\) is defined by
\[F_{0} = 0,\; F_{1} = 1,\; \text{and}\; F_{n} = F_{n-2} + F_{n-1}\text{ for }n\geq 2.\]
Define \(T\in \mathcal{L}(\mathbb{R}^{2})\) by \(T(x,y) = (y,x + y)\).

\begin{enumerate}
    \item[(a)] Show that \(T^{n}(0,1) = (F_{n},F_{n + 1})\) for each nonnegative integer \(n\).
    \item[(b)] Find the eigenvalues of \(T\).
    \item[(c)] Find a basis of \(\mathbb{R}^{2}\) consisting of eigenvectors of \(T\).
    \item[(d)] Use the solution to part (c) to compute \(T^{n}(0,1)\). Conclude that
    \[F_{n} = \frac{1}{\sqrt{5}}\left[\left(\frac{1 + \sqrt{5}}{2}\right)^{n} - \left(\frac{1 - \sqrt{5}}{2}\right)^{n}\right]\]
    for each positive integer \(n\).
    \item[(e)] Use part (d) to conclude that if \(n\) is a nonnegative integer, then the Fibonacci number \(F_{n}\) is the integer that is closest to
    \[\frac{1}{\sqrt{5}}\left(\frac{1 + \sqrt{5}}{2}\right)^{n}.\]
\end{enumerate}
\end{problemBox}

\begin{solutionBox}
\SolutionAnnounce
\doublespacing

\begin{enumerate}
    \item[(a)] 
    \textbf{Base case:} For \(n = 0\), \(T^0(0,1) = (0,1) = (F_0, F_1)\). 
    For \(n = 1\), \(T(0,1) = (1,1) = (F_1, F_2)\).
    
    \textbf{Inductive step:} Assume \(T^n(0,1) = (F_n, F_{n+1})\) for some \(n \geq 1\). Then
    \[
    T^{n+1}(0,1) = T(T^n(0,1)) = T(F_n, F_{n+1}) = (F_{n+1}, F_n + F_{n+1}) = (F_{n+1}, F_{n+2}).
    \]
    Thus by induction, \(T^n(0,1) = (F_n, F_{n+1})\) for all nonnegative integers \(n\).
    
    \item[(b)] 
    We solve \(T(x,y) = \lambda(x,y)\) for eigenvalues \(\lambda\). This gives
    \[(y, x+y) = (\lambda x, \lambda y).\]
    Thus \(y = \lambda x\) and \(x + y = \lambda y\). Then
    \[x + \lambda x = \lambda(\lambda x) \implies x(1 + \lambda) = \lambda^2 x.\]
    If \(x = 0\), then \(y = 0\), which is not an eigenvector. So \(x \neq 0\), giving
    \[1 + \lambda = \lambda^2 \implies \lambda^2 - \lambda - 1 = 0.\]
    Therefore, the eigenvalues are
    \[\lambda_1 = \frac{1 + \sqrt{5}}{2}, \quad \lambda_2 = \frac{1 - \sqrt{5}}{2}.\]
    
    \item[(c)] 
    For \(\lambda_1\), solve \(y = \lambda_1 x\). Taking \(x = 1\) gives the eigenvector \((1, \lambda_1)\). 
    For \(\lambda_2\), similarly we get \((1, \lambda_2)\). 
    Since \(\lambda_1 \neq \lambda_2\) and \(\dim(\R^2) = 2\), these eigenvectors are linearly independent and form a basis. 
    Thus \(B = \{(1, \lambda_1), (1, \lambda_2)\}\) is a basis of \(\R^2\) consisting of eigenvectors of \(T\).
    
    \item[(d)] 
    Let \(B_1 = (1, \lambda_1)\) and \(B_2 = (1, \lambda_2)\). Then
    \[B_1 - B_2 = (0, \lambda_1 - \lambda_2) = (0, \sqrt{5}).\]
    Thus \((0,1) = \frac{1}{\sqrt{5}}(B_1 - B_2)\).
    
    Then
    \[
    T^n(0,1) = \frac{1}{\sqrt{5}}(T^n B_1 - T^n B_2) = \frac{1}{\sqrt{5}}(\lambda_1^n B_1 - \lambda_2^n B_2).
    \]
    Computing the components:
    \[
    T^n(0,1) = \frac{1}{\sqrt{5}}\left(\lambda_1^n - \lambda_2^n, \lambda_1^{n+1} - \lambda_2^{n+1}\right).
    \]
    From part (a), \(T^n(0,1) = (F_n, F_{n+1})\), so comparing the first component gives
    \[F_n = \frac{1}{\sqrt{5}}\left(\lambda_1^n - \lambda_2^n\right) = \frac{1}{\sqrt{5}}\left[\left(\frac{1 + \sqrt{5}}{2}\right)^{n} - \left(\frac{1 - \sqrt{5}}{2}\right)^{n}\right].\]
    
    \item[(e)] 
    From part (d), we have
    \[F_n = \frac{\lambda_1^n}{\sqrt{5}} - \frac{\lambda_2^n}{\sqrt{5}}.\]
    Then
    \[F_n - \frac{\lambda_1^n}{\sqrt{5}} = -\frac{\lambda_2^n}{\sqrt{5}}.\]
    Taking absolute values:
    \[\left|F_n - \frac{\lambda_1^n}{\sqrt{5}}\right| = \left|-\frac{\lambda_2^n}{\sqrt{5}}\right| = \frac{|\lambda_2^n|}{\sqrt{5}}.\]
    Since \(\lambda_2 = \frac{1 - \sqrt{5}}{2} \approx -0.618\), we have \(|\lambda_2| < 1\), so \(|\lambda_2^n| < 1\). Therefore,
    \[\left|F_n - \frac{\lambda_1^n}{\sqrt{5}}\right| < \frac{1}{\sqrt{5}} < \frac{1}{2}.\]
    Thus the distance from \(F_n\), an integer, to \(\frac{\lambda_1^n}{\sqrt{5}}\) is less than \(\frac{1}{2}\), which means \(F_n\) is the integer closest to \(\frac{1}{\sqrt{5}}\left(\frac{1+\sqrt{5}}{2}\right)^n\).
\end{enumerate}

\end{solutionBox}

\end{document}
